\section{Related Work}
\label{sec:related}

\paragraph{Controlled probing.}
Linear classifiers have long been used to track information across network depth \citep{alain2017understanding}. A probe score reflects both the representation and the learner fitted on top of it, motivating matched-capacity probes and control tasks \citep{hewitt2019designing}. The broader probing literature distinguishes information that a readout can recover from information that a model uses for its task \citep{belinkov2022probing}. On synthetic tasks, high probing accuracy can persist for properties that the model provably does not need, directly separating decodability from task relevance \citep{ravichander2021probing}. Causal erasure methods pursue the complementary test: remove linearly identified information and measure the behavioural consequence \citep{ravfogel2020null,elazar2021amnesic}. Our protocol brings these concerns to medical images with patient-disjoint splits, a common 512-dimensional projection, repeated recurring-type controls, and bootstrap intervals. It asks about one verified consumed visual locus rather than selecting the best layer after observing test results.

\paragraph{Representation interventions.}
Activation engineering changes intermediate representations at inference time and measures how outputs respond \citep{turner2023steering,panickssery2024caa}. A probe normal is an appealing intervention direction because it supplies a semantic label and a signed vector from the same representation. Yet a behavioural change under that vector can follow from generic sensitivity to perturbation magnitude, a coordinate-dependent construction, or overlap with another concept direction. Random vectors and a coordinate-permutation sham address the first two possibilities. Fixed clinical probe normals address the third by asking whether the target direction is more effective than plausible semantic alternatives. Our primary statistic recomputes the maximum unrelated clinical effect inside each patient-bootstrap replicate, so the direction-specific comparison is one familywise endpoint rather than five separate pairwise claims.

Intervention families answer different causal questions. Iterative nullspace projection and amnesic probing test whether removing linearly recoverable information damages a downstream task \citep{ravfogel2020null,elazar2021amnesic}. Additive steering instead tests whether a chosen direction can write a change into the forward pass \citep{turner2023steering}. The second test needs an operating scale: an unnormalized vector can appear effective simply because its norm is large, while a fixed Euclidean step can represent very different perturbations across layers or architectures. We use each token's own activation norm as the scale and compare every direction at the same $\alpha$. This choice makes the random, sham, target, and clinical interventions commensurate within a registered cell. It also keeps a generic sensitivity result separate from a claim that the probe normal uniquely represents its named clinical concept.

\paragraph{Medical and multimodal VLM interpretation.}
Depth-wise probing of medical multimodal models has identified losses and gains of classification information across vision towers, connectors, language layers, and semantic outputs \citep{zhu2026lost}. Layer-wise probes of MedCLIP also expose acquisition and device shortcuts in chest radiographs \citep{pedersen2026probes}. These studies motivate tight nuisance and data controls, but they do not intervene on the decoded clinical direction at the same locus. Conversely, controlled visual-concept experiments in lightweight driving VLMs combine probing with probe-vector steering and identify cases where information is linearly present but generation still fails \citep{theodoridis2026probing}. Targeted edits in general VLMs likewise connect linearly available visual features to downstream predictions \citep{rajaram2025lineofsight}. We focus on 7B medical-question answering, preserve the clinical patient split, and compare the target probe normal with equal-norm random, sham, and clinical directions.

Together, this literature motivates three separate questions. Controlled probing asks whether a restricted readout reliably recovers a label. Generic intervention controls ask whether the probe normal moves behaviour beyond non-semantic perturbations. Clinical-direction controls ask whether the response is specific to the named finding. Our experiments align these questions at the same forward-path representation.
